% Tapis de Quillen, with Quillen's offprint and preprint (cote 162-5, typed part): archive cote 162-5, site page 44 (PDF page 45).
% typescript. Draft compared: none.
% AI-transcribed by Claude Opus 5.5 (claude-opus-5-5, adaptive thinking, effort high), Message Batches API, batch t162_full (msgbatch_01UAh3zyWPd454AeDkhnmQ33), 2026-11-13; not yet checked by eye.
\scanpage{162-5/44}
X. De plus, on a la formule \inserted{(du type $g'f. = f'.g'$ ci-dessus, impliquant la formule} de projection habituelle. Ceci posé, les
constructions bien connues en théorie des correspondances permettent
de définir une catégorie $\underline{B}$ dont \unsure{lès|les} objets sont ceux de C, et les
Hom étant donnés par
\[ \mathrm{Hom}_{B}(X,Y) = B(XxY) \quad . \]
On a un isomorphisme canonique $\underline{B} \simeq \underline{B}^{o}$ , qui est l'identité sur les
objets et la symmétrie $B(XxY)\simeq B(YxX)$ sur les flêches, et un foncteur
évident \xout{Bxxx$\underline{B}$x} $F. : C \to \underline{B}$ , qui est l'identité sur les objets et donné
par les graphes sur les flêches. Composé avec la symmétrie précédente,
on trouve un foncteur $F^{\cdot}$: \xout{BX} $C$ \_ $\underline{B}$. On vérifie enfin que le couple
$(F.,F^{\cdot})$ satisfait aux conditions envisagées dans l'alinéa 1 (ceci aussi
devrait être formalisé au cas des correspondances générales ...). \xout{On}
vérifie enfin qu'il a la propriété universelle qu'il faut. Lorsqu'on
se donne un couple $(F!,F'^{\cdot})$ à val eursdans une catégorie B', il faut
trouver une factorisation par B. Elle est toute trouvée sur les
objets, il faut la définir sur les flêches, donc définir, pour X,Y
dans C, une application
\[ (*) \qquad B(XxY) \to \mathrm{Hom}(X',Y') \ . \]
Mais si Z est au\unsure{d}dessus de XxY, d'où $p:Z\to X$ et $q:Z\to Y$, on lui associe
un élément du deuxième membre de (*) comme le composé
\[ X' \overset{F'^{\cdot}(p)}{\hand{$\longrightarrow$}} Z' \overset{F!(q)}{\hand{$\longrightarrow$}} Y' \quad . \]
Il faut prouver que l'élément obtenu en \unsure{remplaçant} Z \inserted{$=Z_o$} par un $Z_1$ cobor\xout{d}
dant au dessus de X \inserted{est le même}. Or si $g: T\to XxY$ réalise le cobordisme, on
voit (regardant seulement T et la décomposition bord(T) = $Z_{o}\amalg Z_1$)
qu'il existe une variété $\overline{T}$ sur $S^1$ (obtenue en ``doublant'' T) \xout{xtxxxx}
\xout{applixatixxxTxxxxx} transversale au dessus d'un couple de poihts $s_o,s_1$
de $S^1$, tels que $Z_i$ soit l'image inverse \inserted{dans $\overline{T}$} de $s_i$ (i=0,1); d'autre part,

\mspage{44}
